\chapter{الإنشاءات الهندسية}\label{ch:g5:geometry}

المسطرة والكوس والفرجار: بهذه الأدوات الثلاث وشيء من
المنهج، يمكن بناء أي شكل موصوف بالكلمات بناءً دقيقًا. ينشئ
هذا الفصل الأعمدة والمتوازيات والمثلثات والرباعيات
الخاصة --- الحرفة أولًا، والنظرية لاحقًا
(\cref{ch:g6:lines,ch:g6:shapes}).

\section{الأعمدة والمتوازيات المارّة بنقطة}

\begin{method}[الإنشاءان الأساسيان]\label{met:g5:geometry:basic}
\emph{العمودي على $d$ المارّ بالنقطة $A$}: أزلق الكوس على
$d$ حتى يبلغ حرفه الآخر ذو \omterm{def:g3:shapes:rightangle}{الزاوية القائمة} النقطة $A$؛ ثم اُرسم؛
وعلّم \omterm{def:g3:shapes:rightangle}{الزاوية القائمة}.

\emph{الموازي للمستقيم $d$ المارّ بالنقطة $A$}: ثبّت مسطرة على
الكوس كسكّة موضوعة على $d$؛ وأزلق الكوس حتى $A$؛
ثم اُرسم. (والتحقّق: المسافة بين المستقيمين هي نفسها في كل
مكان.)
\end{method}

\begin{example}[تسلسل الإنشاءات]\label{ex:g5:geometry:chain}
كثير من الأشكال سلاسل من هاتين الحركتين. فمستطيل $ABCD$ فيه
$AB = 5$ سم و $BC = 3$ سم: اُرسم $[AB]$؛ ثم العمودي عند $A$؛
ثم العمودي عند $B$؛ وعلّم $D$ و $C$ على بُعد $3$ سم من الجهة نفسها؛
ثم صِل. فكل خطوة إنشاء أساسي واحد.
\end{example}

\section{المثلثات بالفرجار}

\begin{method}[مثلث من ثلاثة أضلاع]\label{met:g5:geometry:triangle}
لبناء مثلث أضلاعه $6$ سم و $5$ سم و $4$ سم:
\begin{enumerate}
  \item اُرسم الضلع الأطول، $[AB]$ حيث $AB = 6$ سم؛
  \item وقوسًا مركزه $A$ ونصف قطره $5$ سم؛
  \item وقوسًا مركزه $B$ ونصف قطره $4$ سم؛
  \item وتقاطعهما هو الرأس الثالث $C$؛ ثم صِل.
\end{enumerate}
والفرجار يضمن الأطوال: فكل نقطة من القوس الأول تبعد
$5$ سم بالضبط عن $A$.
\end{method}

\begin{omfigure}
\begin{tikzpicture}[scale=0.95]
  \coordinate (A) at (0,0);
  \coordinate (B) at (6*0.72,0);
  \coordinate (C) at (2.7,2.376);
  \draw[thick] (A) -- (B);
  \draw[thick, omDef] (A) -- (C);
  \draw[thick, omThm] (B) -- (C);
  \draw[omDef, thin] (C) arc[start angle=41.3, delta angle=20,
    radius=3.6];
  \draw[omDef, thin] (C) arc[start angle=41.3, delta angle=-20,
    radius=3.6];
  \draw[omThm, thin] (C) arc[start angle=124.4, delta angle=20,
    radius=2.88];
  \draw[omThm, thin] (C) arc[start angle=124.4, delta angle=-20,
    radius=2.88];
  \fill (A) circle (1.5pt) node[below left, font=\small] {$A$};
  \fill (B) circle (1.5pt) node[below right, font=\small] {$B$};
  \fill (C) circle (1.5pt) node[above, font=\small] {$C$};
  \node[below, font=\small] at (2.16,0) {$6$ سم};
  \node[omDef, font=\small] at (0.85,1.45) {$5$ سم};
  \node[omThm, font=\small] at (3.9,1.35) {$4$ سم};
\end{tikzpicture}

{\small قوسان وتقاطع واحد: المثلث $6$--$5$--$4$ مبنيًّا
بالفرجار. (ولو كان الضلعان القصيران معًا أقصر من
القاعدة، لما التقى القوسان --- والمزيد عن ذلك في
\cref{ch:g7:triangles}.)}
\end{omfigure}

\begin{example}[المتساوي الساقين والمتساوي الأضلاع مجانًا]\label{ex:g5:geometry:special}
فتحة الفرجار نفسها للقوسين $\Rightarrow$ مثلث متساوي الساقين.
وفتحة تساوي القاعدة $\Rightarrow$ متساوي الأضلاع. والفرجار
لا يقيس: بل \emph{ينسخ} الأطوال، وذلك أفضل.
\end{example}

\section{الرباعيات الخاصة بالإنشاء}

\begin{method}[بناء العائلة]\label{met:g5:geometry:quadrilaterals}
\begin{enumerate}
  \item \emph{المستطيل}: زوجان من المتوازيات يتقاطعان في \omterm{def:g3:shapes:rightangle}{زوايا
        قائمة} --- أو سلسلة \cref{ex:g5:geometry:chain}؛
  \item و\emph{المربع}: مستطيل تستعمل أضلاعه الأربعة الطول
        نفسه --- اِنقله بالفرجار؛
  \item و\emph{المعيّن}: أربعة أضلاع متساوية --- اُرسم ضلعًا، ثم
        قوسان من طرفيه بالفتحة نفسها يحددان الرأسين
        الآخرين (فالمعيّن مثلثان متساويا الساقين ملصوقان على
        \omterm{def:g4:geometry:circle}{قطر}).
\end{enumerate}
وبعد البناء، \emph{تحقّق} دائمًا بالأدوات: الكوس على
الأركان، والفرجار على الأضلاع.
\end{method}

\begin{example}[التعرّف من وصف]\label{ex:g5:geometry:recognize}
«\omterm{def:g4:geometry:polygon}{رباعي} أضلاعه الأربعة $4$ سم وفيه \omterm{def:g3:shapes:rightangle}{زاوية قائمة}» ---
اِبنه: فالأضلاع الأربعة المتساوية تفرض شكل معيّن، وضغط
ركن واحد إلى \omterm{def:g3:shapes:rightangle}{زاوية قائمة} يقوّم الأركان الأخرى كلها: فالنتيجة
مربع. وبعض الأوصاف تصف سرًّا شكلًا أخصّ مما
تعلن. (وسبب انتشار \omterm{def:g3:shapes:rightangle}{الزاوية القائمة} إلى الأركان الأربعة
مشروح في \cref{ch:g7:central}.)
\end{example}

\section{الدوائر في الإنشاءات}

\begin{example}[الدائرة أداةً]\label{ex:g5:geometry:circle}
«جِد كل النقط التي تبعد $3$ سم عن $A$»: هي الدائرة التي مركزها $A$
ونصف قطرها $3$ سم. و«جِد كل النقط التي تبعد $3$ سم عن $A$ \emph{و}
$4$ سم عن $B$»: اُرسم الدائرتين؛ وتقاطعاتهما (نقطتان أو نقطة
أو لا نقطة) هي الأجوبة. وخرائط الكنوز وإنشاءات المثلثات
تستعمل هذه الفكرة بالضبط.
\end{example}

\section{تمارين}

\begin{exercise}[$\star$]\label{exo:g5:geometry:1}
اُرسم مستقيمًا $d$ ونقطة $A$ على بُعد $3$ سم منه تقريبًا. وأنشئ
العمودي على $d$ المارّ بالنقطة $A$، ثم الموازي للمستقيم $d$
المارّ بالنقطة $A$. وعلّم \omterm{def:g3:shapes:rightangle}{الزوايا القائمة}.
\end{exercise}

\begin{exercise}[$\star$]\label{exo:g5:geometry:2}
أنشئ مستطيلًا $ABCD$ فيه $AB = 7$ سم و $BC = 4$ سم،
متّبعًا \cref{ex:g5:geometry:chain}. وتحقّق من الزاوية الرابعة
بالكوس: فإن كان الإنشاء دقيقًا كانت قائمة
تلقائيًّا.
\end{exercise}

\begin{exercise}[$\star$]\label{exo:g5:geometry:3}
أنشئ مثلثًا أضلاعه $7$ سم و $6$ سم و $3$ سم
(\cref{met:g5:geometry:triangle}).
\end{exercise}

\begin{exercise}[$\star$]\label{exo:g5:geometry:4}
أنشئ مثلثًا متساوي الساقين قاعدته $5$ سم وساقاه
$7$ سم؛ ثم مثلثًا متساوي الأضلاع طول ضلعه $6$ سم.
\end{exercise}

\begin{exercise}[$\star$]\label{exo:g5:geometry:5}
أنشئ مربعًا طول ضلعه $5$ سم. فأي الأدوات تضمن \omterm{def:g3:shapes:rightangle}{الزوايا
القائمة}؟ وأيّها يضمن تساوي الأضلاع؟
\end{exercise}

\begin{exercise}[$\star$]\label{exo:g5:geometry:6}
أنشئ معيّنًا طول ضلعه $4$ سم وليس مربعًا (اختر بنفسك
فتحة الركن الأول). وتحقّق من الأضلاع الأربعة بالفرجار.
\end{exercise}

\begin{exercise}[$\star$]\label{exo:g5:geometry:7}
اُرسم نقطتين $A$ و $B$ حيث $AB = 5$ سم. وأنشئ كل النقط
التي تبعد $4$ سم عن $A$ و $3$ سم عن $B$. فكم عددها؟
\end{exercise}

\begin{exercise}[$\star$]\label{exo:g5:geometry:8}
اُكتب برنامج إنشاء (خطوات مرقّمة، ونقاط مسمّاة) لمستطيل
$6$ سم $\times$ $2$ سم يعلوه مثلث متساوي الساقين
ساقاه $4$ سم (على هيئة قلم). ودع زميلًا ينفّذه.
\end{exercise}

\begin{exercise}[$\star\star$]\label{exo:g5:geometry:9}
حاول أن تنشئ مثلثًا أضلاعه $8$ سم و $3$ سم و $4$ سم.
فماذا يحدث للقوسين؟ اِشرح في جملة واحدة لماذا لا يمكن أن
يوجد هذا المثلث.
\end{exercise}

\begin{exercise}[$\star\star$]\label{exo:g5:geometry:10}
أنشئ معيّنًا قطراه $6$ سم و $4$ سم،
عالمًا أن قطرَي المعيّن يتقاطعان في منتصفيهما في
\omterm{def:g3:shapes:rightangle}{زاوية قائمة}: اُرسم القطرين \emph{أولًا}، ثم صِل بين أطرافهما
الأربعة.
\end{exercise}

\begin{exercise}[$\star\star$]\label{exo:g5:geometry:11}
مِعزاة مربوطة إلى وتد $A$ بحبل طوله $4$ م، في حقل مستوٍ فيه
سياج مستقيم يمرّ على بُعد $3$ م من $A$. اُرسم مخططًا (1 سم لكل
متر): فأي منطقة تستطيع المعزاة أن تبلغها؟ وأين يحدّها
السياج؟
\end{exercise}
